\cref{itm:ch20-env:directed}.  Since \(\theta\) is semipositive, any
constant no greater than \(\inf\varphi\) is an admissible real model
function.  A rational such constant also proves nonemptiness of the
rational subfamily.  By \Cref{prop:ch20-model-max}, that rational
subfamily is closed under finite maxima and is therefore upward directed.

Now assume \(\{\theta\}\) ample.  Given a real admissible model function
\(\psi\) and \(\varepsilon>0\), use
\Cref{lem:ch20-fixed-form-rational-approximation} to choose a rational
\(\theta\)-psh model function \(q\) with
\(\|q-\psi\|_\infty<\varepsilon\).  Choose a rational number
\(c\) with \(\varepsilon<c<2\varepsilon\).  Then
\[
             \psi-3\varepsilon<q-c\leq\psi\leq\varphi.
\]
Subtracting the rational constant preserves both rationality and
\(\theta\)-pshness.  Taking suprema over rational minorants and then
letting \(\varepsilon\) decrease to zero proves that their supremum
equals the real-model envelope.  The rational minorants, ordered by
pointwise comparison, now supply the asserted increasing net.  No
countability or metrizability is used (\Cref{app:nets}).

The full real-model family itself need not be upward directed, even
under this ample-class hypothesis, by
\Cref{ex:ch20-real-max-irrational-breakpoint}.  The small downward
constant in the preceding argument is what allows the exact obstacle
constraint to survive rational approximation.

